
\documentclass[addpoints]{exam}
\usepackage{amsmath}
\usepackage{amssymb}
\usepackage{color}

% \printanswers

\newcommand{\event}{Lines and Planes in $\mathbb{R}^3$}
\newcommand{\tournament}{}
\def\type{0}

\setlength\fillinlinelength{\textwidth}
\CorrectChoiceEmphasis{\color{red}\bfseries}
\SolutionEmphasis{\color{red}\bfseries}

\setlength\answerclearance{2pt}
\pagestyle{headandfoot}
\runningheadrule
\runningheader{\event}{\tournament}{\ifnum\type=1 Class Set Test \else Full Name:\hspace{1cm} \fi}
\runningfootrule
\runningfooter{}{Page \thepage\ of \numpages}{}

\begin{document}
\begin{center}
    {\huge Grade 12 Calculus \& Vectors \\ \event
    \ifprintanswers \space Key
    \else \space \ifnum\type<2 Test \fi \ifnum\type=1 \textbf{(CLASS SET)} \fi
          \ifnum\type=2 Answer Sheet \fi
    \fi \\ \vspace{3mm}\tournament\vspace{1mm}}
\end{center}

\vspace{.25\textheight}

\begin{center}
    First Name: \ifprintanswers
    \underline{\hspace{3.5cm}}\textbf{KEY}\underline{\hspace{5.5cm}}\\\vspace{5mm}
    \else \underline{\hspace{10cm}}\\ \vspace{5mm} \fi
    Last Name: \underline{\hspace{10cm}}\\ \vspace{5mm}
\end{center}

\noindent\textbf{\underline{Directions:}}
\begin{itemize}
    \ifnum\type=1 \item \textbf{This test is a class set.} Do not write on this paper. \fi
    \item Show all work. State conclusions clearly (intersect / parallel / skew).
    \item Distance formula: $d = \dfrac{|ax_0 + by_0 + cz_0 + d|}{\sqrt{a^2+b^2+c^2}}$.
    \item The test is designed to be completed in 75 minutes.
\end{itemize}
\begin{center}
\textit{For grading use only}\\
\multirowgradetable{1}[pages]
\end{center}

\newpage
\begin{questions}

\section*{Part A --- Multiple Choice (15 marks)}
\vspace{5mm}

\question[1] The vector equation of a line through point $(1,2,3)$ with direction $(4,5,6)$ is:
\begin{choices}
    \choice $(x,y,z) = (4,5,6) + t(1,2,3)$
    \correctchoice $(x,y,z) = (1,2,3) + t(4,5,6)$
    \choice $(x,y,z) = t(1,2,3)(4,5,6)$
    \choice $(x,y,z) = (1,2,3)(4,5,6)$
\end{choices}

\question[1] The parametric equations of the line $\vec{r} = (2,-1,3) + t(1,4,-2)$ are:
\begin{choices}
    \correctchoice $x = 2+t,\; y = -1+4t,\; z = 3-2t$
    \choice $x = 2t,\; y = -t,\; z = 3t$
    \choice $x = 1+2t,\; y = 4-t,\; z = -2+3t$
    \choice $x = t,\; y = 4t,\; z = -2t$
\end{choices}

\question[1] The symmetric equations of a line through $(0,1,-1)$ with direction $(2,-3,1)$
are:
\begin{choices}
    \choice $\dfrac{x}{0} = \dfrac{y}{1} = \dfrac{z}{-1}$
    \correctchoice $\dfrac{x}{2} = \dfrac{y-1}{-3} = \dfrac{z+1}{1}$
    \choice $\dfrac{x-2}{0} = \dfrac{y+3}{1} = \dfrac{z-1}{-1}$
    \choice $2x = -3(y-1) = (z+1)$
\end{choices}

\question[1] The normal vector to the plane $2x - 3y + z = 7$ is:
\begin{choices}
    \choice $(2,-3,7)$
    \choice $(7,-3,1)$
    \correctchoice $(2,-3,1)$
    \choice $(-3,2,1)$
\end{choices}

\question[1] Two lines in $\mathbb{R}^3$ that are not parallel and do not intersect are called:
\begin{choices}
    \choice Collinear
    \correctchoice Skew
    \choice Normal
    \choice Coincident
\end{choices}

\question[1] Two planes are parallel if and only if:
\begin{choices}
    \choice Their equations have the same constant term.
    \correctchoice Their normal vectors are parallel.
    \choice They have the same $x$-intercept.
    \choice They both pass through the origin.
\end{choices}

\question[1] The distance from the origin to the plane $3x - 4y = 10$ is:
\begin{choices}
    \choice 1
    \correctchoice 2
    \choice 5
    \choice 10
\end{choices}

\question[1] Two non-parallel planes in $\mathbb{R}^3$ intersect in:
\begin{choices}
    \choice A point
    \correctchoice A line
    \choice A plane
    \choice They cannot intersect
\end{choices}

\question[1] The direction vector of the line of intersection of two planes can be found by:
\begin{choices}
    \choice Adding their normal vectors
    \correctchoice Taking the cross product of their normal vectors
    \choice Taking the dot product of their normal vectors
    \choice Subtracting their normal vectors
\end{choices}

\question[1] The line $\vec{r} = (1,2,-1) + t(1,-1,2)$ intersects the plane $x+y-z=6$ at:
\begin{choices}
    \choice $(1,2,-1)$
    \correctchoice $(0,3,-3)$
    \choice $(2,1,1)$
    \choice The line is parallel to the plane
\end{choices}

\question[1] Lines with direction vectors $(1,2,-1)$ and $(2,4,-2)$ are:
\begin{choices}
    \choice Perpendicular
    \choice Skew
    \correctchoice Parallel
    \choice Intersecting at a unique point
\end{choices}

\question[1] The distance from $P(2,-1,3)$ to the plane $x + 2y - 2z = 1$ is:
\begin{choices}
    \choice $\dfrac{1}{3}$
    \choice $\dfrac{1}{9}$
    \correctchoice $3$
    \choice $9$
\end{choices}

\question[1] A line is perpendicular to a plane if:
\begin{choices}
    \choice The line's direction vector equals zero
    \correctchoice The line's direction vector is parallel to the plane's normal vector
    \choice The line lies entirely within the plane
    \choice The line's direction vector is perpendicular to the normal
\end{choices}

\question[1] For the plane $ax + by + cz = d$, the vector $(a, b, c)$ is called the:
\begin{choices}
    \choice Direction vector
    \correctchoice Normal vector
    \choice Position vector
    \choice Tangent vector
\end{choices}

\question[1] Three planes can intersect in:
\begin{choices}
    \choice Only a single point
    \choice Only a line
    \correctchoice A point, a line, or have no common intersection
    \choice Always a plane
\end{choices}


\section*{Part B --- Short Answer (33 marks)}

\question Points $A(1,\,-1,\,2)$ and $B(3,\,2,\,-1)$ determine a line $\ell$.
\begin{parts}
    \part[1] Find the direction vector $\vec{d}$ of $\ell$.
    \part[2] Write the \textbf{vector} equation of $\ell$.
    \part[2] Write the \textbf{parametric} equations of $\ell$.
    \part[2] Write the \textbf{symmetric} equations of $\ell$.
    \part[1] Find the coordinates of a third point on $\ell$ (other than $A$ or $B$).
\end{parts}
\begin{solution}[8cm]
\textbf{(a)} $\vec{d} = B - A = (3-1,\,2-(-1),\,-1-2) = \mathbf{(2,\,3,\,-3)}$

\textbf{(b)} $\vec{r} = (1,-1,2) + t(2,3,-3)$

\textbf{(c)} $x = 1+2t,\quad y = -1+3t,\quad z = 2-3t$

\textbf{(d)} $\dfrac{x-1}{2} = \dfrac{y+1}{3} = \dfrac{z-2}{-3}$

\textbf{(e)} Any $t \neq 0, 1$. E.g.\ $t = 2$: $(5,\,5,\,-4)$.
\end{solution}

\question A plane $\pi$ passes through the points $P(1,\,0,\,2)$, $Q(2,\,3,\,1)$,
and $R(-1,\,1,\,0)$.
\begin{parts}
    \part[2] Find two vectors lying in $\pi$.
    \part[2] Find the normal vector $\vec{n}$ to $\pi$ using the cross product.
    \part[2] Write the scalar equation of $\pi$.
    \part[2] Verify that $R$ satisfies the equation. Then determine whether the point
    $S(2,\,2,\,2)$ lies on $\pi$.
\end{parts}
\begin{solution}[9cm]
\textbf{(a)} $\overrightarrow{PQ} = (1,3,-1)$;\quad $\overrightarrow{PR} = (-2,1,-2)$

\textbf{(b)}
\[
\vec{n} = \overrightarrow{PQ}\times\overrightarrow{PR} =
\begin{vmatrix}\vec{i}&\vec{j}&\vec{k}\\1&3&-1\\-2&1&-2\end{vmatrix}
= \vec{i}(-6+1) - \vec{j}(-2-2) + \vec{k}(1+6) = \mathbf{(-5,\,4,\,7)}
\]

\textbf{(c)} Using point $P(1,0,2)$:
\[
-5(x-1)+4(y-0)+7(z-2)=0 \implies \mathbf{-5x+4y+7z = 9}
\]

\textbf{(d)} $R(-1,1,0)$: $-5(-1)+4(1)+7(0) = 5+4 = 9$\checkmark

$S(2,2,2)$: $-5(2)+4(2)+7(2) = -10+8+14 = 12 \neq 9$.\quad $S$ is \textbf{not} on $\pi$.
\end{solution}

\question For each pair of lines, determine whether they \textbf{intersect}, are
\textbf{parallel}, or are \textbf{skew}. If they intersect, find the point.
\begin{parts}
    \part[4] $\ell_1:\;\vec{r} = (1,1,1) + s(2,1,3)$\quad and\quad
              $\ell_2:\;\vec{r} = (3,2,4) + t(1,0,1)$
    \part[4] $\ell_3:\;\vec{r} = (1,0,0) + s(1,1,0)$\quad and\quad
              $\ell_4:\;\vec{r} = (0,0,1) + t(1,0,1)$
\end{parts}
\begin{solution}[10cm]
\textbf{(a)} Set equal component-by-component:
\begin{align*}
1+2s &= 3+t \quad\Rightarrow\quad 2s-t=2 \tag{i}\\
1+s &= 2 \quad\Rightarrow\quad s=1 \tag{ii}\\
1+3s &= 4+t \tag{iii}
\end{align*}
From (ii): $s=1$. From (iii): $4=4+t\Rightarrow t=0$. Check (i): $2-0=2$\checkmark

\textbf{Lines intersect} at $t=0$ (or $s=1$): point $= \mathbf{(3,\,2,\,4)}$.

\textbf{(b)} Direction vectors: $(1,1,0)$ and $(1,0,1)$ are not scalar multiples, so lines are
not parallel. Set equal:
\begin{align*}
1+s &= t \tag{i}\\
s &= 0 \quad\Rightarrow\quad s=0 \tag{ii}\\
0 &= 1+t \quad\Rightarrow\quad t=-1 \tag{iii}
\end{align*}
From (ii): $s=0$; from (iii): $t=-1$; check (i): $1+0=1 \neq -1$.
Contradiction $\Rightarrow$ \textbf{no intersection}. Since not parallel, the lines are
\textbf{skew}.
\end{solution}

\question Distance problems.
\begin{parts}
    \part[3] Find the distance from the point $P(1,\,2,\,3)$ to the plane $2x - y + 2z = 10$.
    \part[3] Find the distance between the parallel planes
    $2x - y + 2z = 10$ and $2x - y + 2z = 1$.
    \part[4] The line $\ell:\; x = 1+2t,\; y = 2-t,\; z = 3+t$ and the plane
    $\pi:\; x + y + z = 12$. Find their point of intersection. Then find the angle between
    $\ell$ and $\pi$.
\end{parts}
\begin{solution}[10cm]
\textbf{(a)} $a=2,\,b=-1,\,c=2,\,d=-10$. $\sqrt{a^2+b^2+c^2}=\sqrt{4+1+4}=3$.
\[
d = \frac{|2(1)-1(2)+2(3)-10|}{3} = \frac{|2-2+6-10|}{3} = \frac{4}{3}
\approx \mathbf{1.33 \text{ units}}
\]

\textbf{(b)} Pick any point on $2x-y+2z=1$, e.g.\ $(0,0,1/2)$:
\[
d = \frac{|2(0)-0+2(1/2)-10|}{3} = \frac{|1-10|}{3} = \frac{9}{3} = \mathbf{3 \text{ units}}
\]
(Equivalently: $d = |10-1|/\sqrt{4+1+4} = 9/3 = 3$.)

\textbf{(c)} Substitute $\ell$ into $\pi$:
$(1+2t)+(2-t)+(3+t) = 12 \implies 6+2t=12 \implies t=3$.
Point: $(1+6,\,2-3,\,3+3) = \mathbf{(7,\,-1,\,6)}$.

Angle between line and plane: $\sin\phi = \dfrac{|\vec{d}\cdot\vec{n}|}{|\vec{d}||\vec{n}|}$
where $\vec{d}=(2,-1,1)$ and $\vec{n}=(1,1,1)$.
\[
\vec{d}\cdot\vec{n} = 2-1+1 = 2;\quad |\vec{d}|=\sqrt{6};\quad |\vec{n}|=\sqrt{3}
\]
\[
\sin\phi = \frac{2}{\sqrt{6}\cdot\sqrt{3}} = \frac{2}{3\sqrt{2}} = \frac{\sqrt{2}}{3}
\implies \phi = \arcsin\!\left(\tfrac{\sqrt{2}}{3}\right) \approx \mathbf{28.1^\circ}
\]
\end{solution}

\end{questions}
\end{document}
